\bookchapter{บทที่ 1: สวัสดี คอมพิวเตอร์ส่วนตัวของฉัน}

\chapterintro{โทรศัพท์ของคุณ จริงๆ แล้วคือซูเปอร์คอมพิวเตอร์ส่วนตัวในกระเป๋าของคุณ มันทรงพลังกว่าคอมพิวเตอร์ที่เคยส่งมนุษย์ไปดวงจันทร์นับไม่ถ้วน แต่กลับนอนนิ่งอยู่ในฝ่ามือของคุณ วันนี้ มาทำความรู้จัก “ยักษ์ใจเงียบ” ตัวนี้กัน}

\booksection{“อวัยวะภายใน” ของฮาร์ดแวร์โทรศัพท์: วิเคราะห์เชิงลึก}

โทรศัพท์เล็กแต่มีอวัยวะครบ มาทำความรู้จักองค์ประกอบที่แม่นยำเหล่านี้ทีละชิ้นเหมือนในคาบเรียนกายวิภาค

\booksubsection{SoC (System on Chip): “สมอง+หัวใจ” ของโทรศัพท์}

SoC ย่อมาจาก System on Chip แปลได้ว่า “บรรจุทั้งระบบไว้ในชิปเดียว” มันไม่ใช่แค่ CPU (หน่วยประมวลผลกลาง) แต่รวม CPU, GPU (หน่วยประมวลผลกราฟิก), NPU (หน่วยประมวลผลโครงข่ายประสาทเทียม), ISP (หน่วยประมวลผลสัญญาณภาพ), เบสแบนด์ (รับผิดชอบการสื่อสาร) และอื่นๆ ไว้ทั้งหมด

เหมือนการสร้างอาคารรัฐบาล โรงไฟฟ้า ศูนย์สื่อสาร สถานีตำรวจของเมืองไว้ในอาคารเดียว มีประสิทธิภาพสูงมาก

\booksubsection{CPU: แก่นกลางของการคำนวณทั่วไป}

CPU เป็นองค์ประกอบหลักของ SoC รับผิดชอบการคำนวณส่วนใหญ่ CPU ของโทรศัพท์ใช้สถาปัตยกรรม ARM ซึ่งต่างจากสถาปัตยกรรม x86 ของคอมพิวเตอร์ สถาปัตยกรรม ARM มีลักษณะเด่นที่ชุดคำสั่งแบบย่อ (RISC) ใช้พลังงานต่ำกว่า เหมาะกับอุปกรณ์พกพามากกว่า

\textbf{การออกแบบ CPU แบบหลักใหญ่-หลักเล็ก}: CPU ของโทรศัพท์สมัยใหม่ใช้สถาปัตยกรรม “หลักใหญ่-หลักเล็ก” (เทคโนโลยี big.LITTLE ของ ARM) โดยแบ่งหลักออกเป็นระดับประสิทธิภาพต่างกัน:

\begin{itemize}
  \item \textbf{ซูเปอร์คอร์/บิ๊กคอร์}: ประสิทธิภาพสูงสุด รับผิดชอบงานหนัก เช่น โหลดเกม เปิดแอป คำนวณซับซ้อน ความถี่สูงสุด (3GHz+) แต่ใช้พลังงานมากที่สุด
  \item \textbf{คอร์กลาง}: สมดุลระหว่างประสิทธิภาพและพลังงาน จัดการงานโหลดปานกลางประจำวัน เช่น เลื่อนดูโซเชียล ดูวิดีโอ สลับแอป
  \item \textbf{คอร์เล็ก/คอร์ประหยัดพลังงาน}: ใช้พลังงานต่ำสุด จัดการงานเบื้องหลังเบาๆ เช่น สแตนด์บาย ฟังเพลง รับข้อความ แทบไม่ใช้ไฟ
\end{itemize}

หลักสามระดับทำงานร่วมกัน: ตัวกำหนดเวลาระบบกระจายงานอัตโนมัติตามโหลด — งานเบาให้คอร์เล็กเพื่อประหยัดไฟ งานหนักให้คอร์ใหญ่ทำเต็มที่ เสร็จแล้วสลับกลับคอร์เล็ก นี่คือเหตุผลที่โทรศัพท์เล่นเกมลื่นได้และสแตนด์บายได้ทั้งวัน

\textbf{ลักษณะเด่นของสถาปัตยกรรมแต่ละแบรนด์}:

\begin{itemize}
  \item \textbf{Qualcomm Snapdragon 8 Elite}: ใช้สถาปัตยกรรม CPU Oryon ที่พัฒนาขึ้นเอง (จากการเข้าซื้อ Nuvia) ประกอบด้วยซูเปอร์คอร์ 2 ตัว (4.32GHz) + คอร์ประสิทธิภาพ 6 ตัว (3.53GHz) ทั้งหมดผลิตด้วยกระบวนการ TSMC 3nm รุ่นที่ 2 ละทิ้งสถาปัตยกรรม ARM ทั่วไปอย่างสิ้นเชิง ประสิทธิภาพ single-core และ multi-core พัฒนาแบบก้าวกระโดด GPU เป็น Adreno 830 ใช้สถาปัตยกรรมแบบ sliced เพิ่มประสิทธิภาพกราฟิก 40\% NPU Hexagon มีความสามารถในการประมวลผลประมาณ 45 TOPS เพิ่มประสิทธิภาพการอนุมาน AI 45\% (หมายเหตุ: รุ่นพิเศษ Samsung Galaxy มีความถี่สูงกว่า ซูเปอร์คอร์สูงสุด 4.47GHz)
  \item \textbf{MediaTek Dimensity 9400}: ซูเปอร์คอร์ 1 ตัว (Cortex-X925, 3.62GHz) + คอร์ประสิทธิภาพ 3 ตัว (Cortex-X4, 3.3GHz) + คอร์ประหยัดพลังงาน 4 ตัว (Cortex-A720, 2.4GHz) สถาปัตยกรรม All-Big-Core ต่อเนื่อง กระบวนการ TSMC 3nm รุ่นที่ 2 GPU เป็น Immortalis-G925 (12 คอร์) เพิ่มประสิทธิภาพสูงสุด 41\% และ Ray Tracing 40\% หน่วยประมวลผล AI เป็น NPU 890 รุ่นที่ 8 รองรับการรัน Large Model บนอุปกรณ์ การฝึก LoRA บนอุปกรณ์ และการสร้างวิดีโอบนอุปกรณ์
  \item \textbf{Apple A18 Pro}: คอร์ประสิทธิภาพ 2 ตัว (\textasciitilde{}4.04GHz) + คอร์ประหยัดพลังงาน 4 ตัว (\textasciitilde{}2.4GHz) กระบวนการ TSMC 3nm รุ่นที่ 2 GPU 6 คอร์ รองรับ Ray Tracing ระดับฮาร์ดแวร์ NPU 16 คอร์ (ประมาณ 35 TOPS) มอบประสิทธิภาพการประมวลผลที่แข็งแกร่งสำหรับ AI บนอุปกรณ์และ Apple Intelligence สถาปัตยกรรมที่พัฒนาเองของ Apple ยังคงเป็นผู้นำด้านประสิทธิภาพ single-core และประสิทธิภาพการใช้พลังงาน สถาปัตยกรรม Unified Memory ให้ข้อได้เปรียบแบนด์วิดท์สูงสำหรับการอนุมาน AI
  \item \textbf{Huawei Kirin 9020}: ติดตั้งใน Mate 70 Series ใช้สถาปัตยกรรม CPU Taishan ที่พัฒนาขึ้นเอง ความสามารถของ NPU และ ISP ได้รับการปรับปรุงอย่างต่อเนื่อง รองรับ 5.5G และการสื่อสารผ่านดาวเทียม เป็นตัวแทนความก้าวหน้าล่าสุดของชิปที่ผลิตในประเทศ
\end{itemize}

\booksubsection{GPU: เอนจินเรนเดอร์กราฟิก}

GPU (Graphics Processing Unit) รับผิดชอบกราฟิกและงานคำนวณแบบขนานโดยเฉพาะ จำนวนคอร์มากกว่า CPU มาก (มักจะหลายร้อยถึงหลายพันคอร์) แต่ละคอร์ทำงานได้ค่อนข้างเรียบง่าย เหมาะกับการประมวลผลพร้อมกันจำนวนมาก

\textbf{ขั้นตอนการเรนเดอร์กราฟิก}:

\begin{enumerate}
  \item \textbf{ประมวลผลเวอร์เท็กซ์}: คำนวณพิกัดตำแหน่งของแต่ละเวอร์เท็กซ์ในโมเดล 3D
  \item \textbf{Rasterization}: แปลงกราฟิก 3D เป็นพิกเซลบนหน้าจอ 2D
  \item \textbf{Pixel Shading}: คำนวณสี พื้นผิว และเอฟเฟกต์แสงสำหรับแต่ละพิกเซล
  \item \textbf{การสร้างภาพขั้นสุดท้าย}: รวมเลเยอร์ทั้งหมดเป็นภาพสุดท้าย ส่งไปแสดงบนหน้าจอ
\end{enumerate}

\textbf{ผลกระทบของ GPU ต่อการใช้งานประจำวัน}: ไม่ใช่แค่เกม! การเลื่อนอินเทอร์เฟซ เอฟเฟกต์แอนิเมชัน การเล่นวิดีโอ ฟิลเตอร์รูปภาพ แอป AR ล้วนพึ่งพา GPU โทรศัพท์ที่ GPU อ่อนแอ การเลื่อนอินเทอร์เฟซอาจกระตุก เกมยิ่งค้างโดยตรง

\textbf{GPU แต่ละแบรนด์}: Qualcomm ใช้ Adreno GPU ที่พัฒนาขึ้นเอง; MediaTek ใช้ ARM Mali GPU; Apple ใช้ GPU ที่พัฒนาขึ้นเอง ทั้งประสิทธิภาพและประสิทธิภาพการใช้พลังงานล้วนเป็นมาตรฐานอุตสาหกรรม

\booksubsection{NPU: เอนจินเร่ง AI}

NPU (Neural Processing Unit) เป็นโปรเซสเซอร์ที่ออกแบบมาเฉพาะสำหรับการคำนวณโครงข่ายประสาทเทียม เมื่อเทียบกับ CPU และ GPU แล้ว NPU มีประสิทธิภาพสูงกว่าและใช้พลังงานต่ำกว่าในการประมวลผลงาน AI

\textbf{หลักการคำนวณของโครงข่ายประสาทเทียม}: การคำนวณหลักของ NPU คือการคูณเมทริกซ์และการคอนโวลูชัน — ซึ่งเป็นการคำนวณพื้นฐานและเกิดขึ้นบ่อยที่สุดในโครงข่ายประสาทเทียม NPU ใช้การเพิ่มประสิทธิภาพระดับฮาร์ดแวร์ สามารถประมวลผลการคูณเมทริกซ์ปริมาณมากแบบขนาน เร็วกว่า CPU หลายสิบเท่า แต่ใช้พลังงานน้อยกว่ามาก

\textbf{ขั้นตอนการประมวลผล AI}:

\begin{enumerate}
  \item \textbf{การป้อนข้อมูล}: ภาพจากกล้อง เสียงจากไมโครโฟน ข้อมูลเซ็นเซอร์ ฯลฯ
  \item \textbf{การประมวลผลล่วงหน้า}: ปรับรูปแบบและขนาดข้อมูลให้เหมาะสมกับความต้องการอินพุตของโครงข่ายประสาทเทียม
  \item \textbf{การคำนวณอนุมาน}: NPU โหลดโมเดลที่ฝึกเสร็จแล้ว ทำการคำนวณ forward propagation กับข้อมูลอินพุต และแสดงผลลัพธ์
  \item \textbf{การประมวลผลภายหลัง}: แปลงผลลัพธ์ AI เป็นรูปแบบที่ใช้งานได้ เช่น พิกัดใบหน้า เนื้อหาข้อความ แท็กสถานการณ์
\end{enumerate}

\textbf{สถานการณ์การใช้งานจริง}:

\begin{itemize}
  \item \textbf{การถ่ายภาพเชิงคำนวณ}: การถ่ายภาพกลางคืนแบบรวมเฟรม โหมดหน้าชัดหลังเบลอ การประมวลผล HDR การปรับแต่งใบหน้าด้วย AI
  \item \textbf{ผู้ช่วยเสียง}: การรู้จำเสียง ความเข้าใจภาษาธรรมชาติ การสังเคราะห์เสียง
  \item \textbf{การแปลแบบเรียลไทม์}: คำบรรยาย AI การแปลบทสนทนา
  \item \textbf{การรู้จำอัจฉริยะ}: การรู้จำอักขระ (OCR) การรู้จำวัตถุ การรู้จำฉาก
  \item \textbf{Large Model บนอุปกรณ์}: ชิปเรือธงรุ่นล่าสุดรองรับการรัน AI Large Model ขนาดเล็ก (ระดับ 7B พารามิเตอร์) บนโทรศัพท์ ทำให้สามารถสนทนา AI แบบออฟไลน์ สรุปเอกสาร เขียนข้อความอัจฉริยะ และฟังก์ชันอื่นๆ ได้ Snapdragon 8 Elite มี NPU Hexagon ประมวลผลประมาณ 45 TOPS Dimensity 9400 มี NPU 890 ที่รองรับ Large Model บนอุปกรณ์เช่นกัน (เพิ่มความเร็วการอนุมาน LLM 80\%) Apple A18 Pro มี NPU 16 คอร์ ให้พลังประมวลผลประมาณ 35 TOPS
  \item \textbf{การประมวลผลภาพ AI}: AI Eraser (ลบคนเดินผ่านในรูปด้วยคลิกเดียว) AI Expand (ขยายขอบภาพ) AI Style Transfer การสร้างภาพจากข้อความ ฯลฯ
\end{itemize}

\booksubsection{ISP: โปรเซสเซอร์สัญญาณภาพ}

ISP ประมวลผลข้อมูลดิบจากกล้อง แปลงสัญญาณแสงเป็นภาพถ่ายที่สวยงาม

\textbf{ประสิทธิภาพอย่างไรดู?}ดูรุ่น! Apple A18 Pro, Qualcomm Snapdragon 8 Elite, MediaTek Dimensity 9400, Huawei Kirin 9020 ล้วนเป็นระดับฟลากชิป ระดับกลางให้ดู Snapdragon 8s Gen3/7+ Gen3 และ Dimensity 8400 อย่าดูแค่จำนวนคอร์ สถาปัตยกรรมและกระบวนการผลิต (จำนวนนาโนเมตร) สำคัญกว่า 3nm เป็นกระบวนการผลิตที่ล้ำสุดตอนนี้ สมรรถนะต่อพลังงานดีกว่า 5nm และ 4nm อย่างชัดเจน

\booksubsection{RAM: โต๊ะทำงานชั่วคราว}

RAM เหมือนโต๊ะทำงานของคุณ ยิ่งใหญ่เท่าไร ยิ่งเปิดไฟล์พร้อมกันได้มากขึ้น 12GB เป็นมาตรฐานปัจจุบัน 16GB เหมาะกับผู้ใช้งานหนัก 24GB ขึ้นไปเป็นทางเลือกของเกมเมอร์และกีก

ชนิดของ RAM ก็สำคัญเช่นกัน: LPDDR5 เร็วกว่าและประหยัดไฟกว่า LPDDR4X แอปเปิลแม้ RAM จะไม่ใหญ่ แต่ปรับแต่งระบบได้ดี จึงดู “พอใช้”

\textbf{ถ้า RAM ไม่พอจะเกิดอะไรขึ้น?}แอปจะถูกระบบฆ่าออกจากพื้นหลัง พอสลับกลับมาก็ต้องโหลดใหม่ทั้งหมด นี่คือเหตุผลที่บางทีสลับออกจากเกมไปแชท WeChat สักครู่ กลับมาแล้วเกมหายไปเลย

\booksubsection{หน่วยเก็บข้อมูล (ROM/Flash): คลังสินค้าระยะยาว}

หน่วยเก็บข้อมูลเหมือนตู้หนังสือของคุณ เก็บรูป วิดีโอ แอปทั้งหมด ปัจจุบันมาตรฐานคือ UFS 3.1 หรือ UFS 4.0 เร็วกว่า eMMC แบบเก่าหลายเท่า

\textbf{ทำไมความเร็วของหน่วยเก็บข้อมูลถึงสำคัญ?}มันกำหนดความเร็วในการเปิดแอป การโหลดเกม และการบันทึกรูป ใช้ไปนานๆ หน่วยเก็บข้อมูลที่ช้าจะทำให้โทรศัพท์ค้างมากขึ้นเรื่อยๆ

\booksubsection{หน้าจอ: ดวงตาของคุณ}

หน้าจอเป็นหน้าต่างหลักในการโต้ตอบระหว่างคุณกับโทรศัพท์ คุณภาพดีหรือไม่มีผลโดยตรงต่อประสบการณ์การใช้งาน

\begin{itemize}
  \item \textbf{วัสดุ}: OLED สีสันสดใส คอนทราสต์สูง ประหยัดไฟกว่า (ตอนแสดงสีดำพิกเซลจะไม่เปล่งแสง) แต่ความสว่างต่ำอาจมีอาการกระพริบที่ทำร้ายตา ส่วน LCD ดูแลตากว่า ต้นทุนต่ำ แต่คอนทราสต์และสีสันไม่เท่า OLED
  \item \textbf{ความละเอียด}: 1080P เป็นกระแสหลัก 2K คมชัดกว่าแต่กินไฟกว่า แนวคิด "จอเรตินา" ของแอปเปิลก็คือ: ในระยะดูที่เหมาะสม ตาคนจะแยกพิกเซลแต่ละจุดไม่ออก
  \item \textbf{อัตรารีเฟรช}: 60Hz คือมาตรฐานพื้นฐาน 90Hz/120Hz ลื่นกว่า ตอนปัดหน้าจอจะไหลลื่นเหมือนไหม อัตรารีเฟรชสูงกินไฟมากขึ้น หลายโทรศัพท์จึงรองรับอัตรารีเฟรชแบบปรับตัว (LTPO) ปรับตามเนื้อหาที่แสดง
  \item \textbf{อัตราการสุ่มสัญญาณหน้าจอสัมผัส}: มีผลต่อความเร็วการตอบสนองของหน้าจอสัมผัส สายเกมควรดูจุดนี้ ยิ่งสูงยิ่งตอบสนองตามมือ
  \item \textbf{วิธีปรับแสง}: แบบ DC ดูแลตากว่าแบบ PWM ตอนความสว่างต่ำจะไม่เหนื่อยตาง่าย ตอนนี้หลายโทรศัพท์รองรับ PWM ความถี่สูง ก็เพื่อดูแลตาเช่นกัน
\end{itemize}

\booksubsection{เทคโนโลยี HDR ของหน้าจอ}

HDR (High Dynamic Range) ทำให้หน้าจอแสดงส่วนสว่างที่สว่างขึ้นและส่วนมืดที่มืดลงได้พร้อมกัน ระดับความลึกของภาพดีขึ้นอย่างมาก

\begin{itemize}
  \item \textbf{หลักการทำงาน}: หน้าจอทั่วไปมีช่วงความสว่างจำกัด (SDR) ช่องว่างระหว่างจุดสว่างสุดกับมืดสุดไม่พอใหญ่ หน้าจอ HDR มีความสว่างพีคสูงกว่า (มัก 1,000 นิตขึ้นไป) คอนทราสต์ใหญ่กว่า ช่วงสีกว้างกว่า แสดงรายละเอียดได้มากกว่า
  \item \textbf{กระแสหลักมาตรฐาน}: HDR10 (พื้นฐานที่สุด ฟรีและโอเพนซอร์ส), HDR10+ (ผลักดันโดยซัมซุง ใช้ข้อมูลเมตาแบบไดนามิก), Dolby Vision (ดอลบีวิชัน ภาพดีที่สุดแต่ต้องจ่ายค่าลิขสิทธิ์) โทรศัพท์จะรองรับมาตรฐานไหนขึ้นอยู่กับฮาร์ดแวร์หน้าจอและลิขสิทธิ์ของระบบ
  \item \textbf{ผลลัพธ์จริง}: ตอนดูวิดีโอ HDR แสงแดดจะจ้าขึ้น เงามืดจะลึกขึ้น สีจะอุดมขึ้น แต่เงื่อนไขคือวิดีโอนั้นต้องผลิตมาในรูปแบบ HDR และหน้าจอก็ต้องรองรับการแสดง HDR ด้วย
\end{itemize}

\booksubsection{เทคโนโลยี HDR สำหรับภาพถ่าย}

HDR ตอนถ่ายรูปกับ HDR ของหน้าจอเป็นคนละเรื่องกัน ภาพถ่าย HDR คือการใช้อัลกอริทึมแก้ปัญหา “สว่างเกินไปจนขาวโพลน มืดเกินไปจนดำสนิท”

\begin{itemize}
  \item \textbf{หลักการของอัลกอริทึม}: ตอนกดชัตเตอร์ โทรศัพท์จริงๆ ถ่ายหลายภาพที่ค่าแสงต่างกัน (ภาพขาดแสงเก็บรายละเอียดส่วนสว่าง ภาพเกินแสงเก็บรายละเอียดส่วนมืด ภาพปกติเก็บโทนกลาง) แล้วนำมารวมกันเป็นภาพเดียวที่มีช่วงไดนามิกสูง
  \item \textbf{ขั้นตอนการประมวลผล}: จับภาพหลายเฟรม จัดตำแหน่งและชดเชย (กันมือสั่น) ผสานค่าแสง/ปรับโทนสี แก้สีให้ถูกต้อง ส่งออกภาพสุดท้าย
  \item \textbf{การปรับปรุงผลลัพธ์}: ตอนนี้โทรศัพท์ทำ "HDR แบบเรียลไทม์" ได้แล้ว ภาพที่เห็นในช่องมองภาพคือผลหลังประมวลผล HDR แบบเห็นอย่างไรได้อย่างนั้น Smart HDR ของแอปเปิลและ HDR อัตโนมัติของแอนดรอยด์ก็ทำได้ดีแล้วทั้งคู่
\end{itemize}

\booksubsection{แบตเตอรี่และการชาร์จ: แหล่งพลังงาน}

แบตเตอรี่เป็นชิ้นส่วนที่ “อายุสั้น” ที่สุดของโทรศัพท์ ใช้ไป 2-3 ปีจะเสื่อมสภาพอย่างเห็นได้ชัด ปัจจุบันความจุมาตรฐานคือ 5500-7500mAh

\begin{itemize}
  \item \textbf{กำลังชาร์จ}: ค่ายแอนดรอยด์แข่งกันถึงขั้น 100W หรือแม้ 200W ชาร์จเต็มในสิบกว่านาที ฝั่งแอปเปิลยังวนเวียนอยู่ที่ 20-30W ชาร์จเร็วทำร้ายแบตไหม? ถ้าใช้อุปกรณ์แท้ก็ไม่เป็นปัญหาใหญ่ โทรศัพท์มีระบบควบคุมความร้อนอยู่แล้ว
  \item \textbf{การชาร์จไร้สาย}: สะดวกแต่ประสิทธิภาพต่ำ ร้อนสะสมเยอะ เหมาะกับการวางชาร์จเอาไว้ในออฟฟิศ ไม่เหมาะกับตอนรีบออกจากบ้าน
  \item \textbf{การชาร์จย้อนกลับ}: โทรศัพท์กลายเป็นพาวเวอร์แบงก์ ชาร์จให้หูฟัง นาฬิกา ใช้ในยามฉุกเฉิน
\end{itemize}

\booksubsection{กล้อง: เวทมนตร์แห่งแสงและเงา}

ปัจจุบันกล้องโทรศัพท์มักมี 3-4 ตัว แต่ละตัวมีหน้าที่ต่างกัน:

\begin{itemize}
  \item \textbf{กล้องหลัก}: ใช้บ่อยที่สุด เซ็นเซอร์ใหญ่ที่สุด คุณภาพภาพดีที่สุด
  \item \textbf{อัลตร้าไวด์}: ถ่ายวิว อาคาร มองเห็นได้กว้างกว่า
  \item \textbf{เทเลโฟโต}: ถ่ายของไกลๆ มีทั้งซูมออปติคอลและซูมดิจิทัลให้แยกกัน
  \item \textbf{มาโคร}: ถ่ายสิ่งเล็กๆ ในระยะใกล้มาก
\end{itemize}

\textbf{จำนวนพิกเซลไม่ใช่ทุกอย่าง!}ขนาดเซ็นเซอร์ (เซ็นเซอร์ใหญ่กว่าก็กดทับได้หมด) ขนาดรูรับแสง (f/1.8 ใหญ่กว่า f/2.4) และระบบกันสั่น (OIS สำคัญมาก) ล้วนมีผลต่อคุณภาพภาพ ส่วนการถ่ายภาพเชิงคำนวณ (อัลกอริทึม AI) สำคัญขึ้นเรื่อยๆ ช่วยรวมหลายภาพให้เป็นภาพเดียวที่ดีได้

\booksubsection{เลนส์เทเลโฟโตแบบเปอริสโคป}

เลนส์เทเลโฟโตแบบเดิมต้องมีเส้นทางแสงยาว ทำให้โทรศัพท์หนาขึ้นมาก แบบเปอริสโคปใช้ปริซึมมุม 90 องศาหักเหแสง ให้เส้นทางแสงวางแนวนอนไปตามตัวเครื่อง จึงคงความบางเบาไว้ได้พร้อมกับมีซูมออปติคอล 5 เท่าหรือแม้ 10 เท่า

\begin{itemize}
  \item \textbf{หลักการโครงสร้าง}: แสงเข้าจากเลนส์ ปริซึมหักเห 90 องศา เดินตามความกว้างของตัวเครื่องผ่านกลุ่มเลนส์ ถึงเซ็นเซอร์ เท่ากับเอากล้องส่องทางไกลที่ตั้งตรงอยู่มาวางให้นอนลง
  \item \textbf{คุณสมบัติทางทัศนศาสตร์}: อัตราซูมออปติคอลสูง (5x-10x) คุณภาพภาพดีกว่าซูมดิจิทัลอย่างชัดเจน แต่ขนาดเซ็นเซอร์มักเล็กกว่ากล้องหลัก ในที่แสงน้อยจึงทำได้เพียงระดับทั่วไป
  \item \textbf{จุดแข็งด้านภาพ}: ถ่ายดวงจันทร์ ถ่ายอาคารไกลๆ ถ่ายใกล้ชิดนักแสดงบนเวที สิ่งที่โทรศัพท์ทั่วไปเอาไม่ถึง
\end{itemize}

\booksubsection{OIS กันสั่นแบบออปติคอล vs EIS กันสั่นแบบอิเล็กทรอนิกส์}

ระบบกันสั่นคือเทคโนโลยีสำคัญที่มีผลต่อคุณภาพภาพถ่ายและวิดีโอ สองวิธีนี้ทำงานต่างกันคนละแบบ:

\begin{itemize}
  \item \textbf{OIS (กันสั่นแบบออปติคอล)}: ย้ายเลนส์หรือเซ็นเซอร์ด้วยกลไกทางกายภาพเพื่อหักล้างอาการมือสั่น ในเลนส์มีไจโรสโคปขนาดจิ๋วตรวจจับทิศทางและขนาดของการสั่น แล้วสั่งมอเตอร์ขยับกลุ่มเลนส์ไปทางตรงข้าม ข้อดีคือไม่เสียคุณภาพภาพ ข้อเสียคือต้นทุนสูงและกินพื้นที่ภายใน
  \item \textbf{EIS (กันสั่นแบบอิเล็กทรอนิกส์)}: ใช้อัลกอริทึมซอฟต์แวร์ตัดและเลื่อนภาพเพื่อหักล้างการสั่น ตอนอัดวิดีโอโทรศัพท์จะเก็บภาพที่กว้างกว่าปกติเล็กน้อย แล้วตัดตามข้อมูลไจโรสโคปให้เหลือเฉพาะส่วนที่นิ่ง ข้อดีคือไม่ต้องมีฮาร์ดแวร์เพิ่ม ข้อเสียคือภาพถูกตัด (มุมมองแคบลง) และในที่แสงน้อยอาจเบลอ
  \item \textbf{การเลือกในทางปฏิบัติ}: ถ่ายรูปให้ดู OIS ก่อน ส่วนอัดวิดีโอใช้ทั้งสองอย่างร่วมกันจะได้ผลดีที่สุด ตอนนี้โทรศัพท์ฟลากชิปจำนวนมากรองรับเซ็นเซอร์กันสั่นแบบเลื่อนตำแหน่ง (Sensor-Shift) ซึ่งได้ผลดีกว่ากันสั่นที่เลนส์แบบเดิม
\end{itemize}

\booksubsection{ซูมออปติคอล vs ซูมดิจิทัล}

\begin{itemize}
  \item \textbf{ซูมแบบออปติคอล}: ขยับกลุ่มเลนส์ด้วยกลไกทางกายภาพเพื่อเปลี่ยนระยะโฟกัส คุณภาพภาพไม่เสียหาย ซูมออปติคอลบนโทรศัพท์มักมีเฉพาะอัตราที่กำหนดไว้ (เช่น 2x, 3x, 5x, 10x) ไม่สามารถซูมต่อเนื่องได้เหมือนกล้อง
  \item \textbf{ซูมดิจิทัล}: ขยายส่วนกลางของภาพให้ใหญ่ขึ้น แก่นแท้คือการตัดภาพ + ขยายแบบแทรกค่า ยิ่งขยายมาก คุณภาพยิ่งแย่ ซูมดิจิทัล 100 เท่าดูเท่หรู แต่ภาพจริงๆ เบลอมากแล้ว
  \item \textbf{ซูมแบบผสม}: ผสมซูมออปติคอลกับดิจิทัลเข้าด้วยกัน ในช่วงระหว่างอัตราซูมออปติคอลจะใช้อัลกอริทึมรวมภาพหลายเฟรมชดเชยคุณภาพที่เสียไป ได้ผลดีกว่าซูมดิจิทัลล้วน แต่ยังสู้ซูมออปติคอลไม่ได้
\end{itemize}

\booksubsection{เสียง: ลำโพงและชิปแสดงผลเฉพาะทาง}

\textbf{ลำโพงเดี่ยว vs ลำโพงคู่}:

\begin{itemize}
  \item \textbf{ลำโพงเดียว}: มีหน่วยส่งเสียงเดียวด้านล่าง เสียงเอียงไปข้างเดียว ไม่มีความรู้สึกสเตอริโอ พบได้ทั่วไปในโทรศัพท์ระดับกลาง-ล่าง
  \item \textbf{ลำโพงคู่ (เสียงสเตอริโอ)}: ลำโพงแยกกันบนล่างข้างละตัว สร้างสนามเสียงสเตอริโอ ตอนดูวิดีโอฟังเพลงจะแยกซ้ายขวาชัด ให้ความเพลิดเพลินดีกว่าหลายเท่า การให้ตัวเรียกเสียงทำหน้าที่ลำโพงบนด้วยเป็นแนวทางที่พบบ่อย
  \item \textbf{ความต่างด้านคุณภาพเสียง}: ลำโพงคู่ไม่ได้แค่ "เสียงดังขึ้น" สำคัญที่ความรู้สึกเป็นพื้นที่และมีมิติ ลองใช้ลำโพงเดียวดูหนังใหญ่ๆ จะรู้สึกเหมือนเหลืออะไรบางอย่าง
\end{itemize}

\textbf{เที่ยวเล่นโทรศัพท์มือถือโดดเดี่ยวแสดงชิปแผ่น}:

โทรศัพท์เกมมิ่งและโทรศัพท์สายเพอร์ฟอร์แมนซ์บางรุ่นมีชิปประมวลผลภาพแยกต่างหาก หน้าที่ของมันคือ:

\begin{itemize}
  \item \textbf{แทรกเฟรม}: ยกภาพเกม 30fps ขึ้นเป็น 60fps หรือ 120fps ด้วยการแทรกเฟรม ภาพลื่นขึ้น ขั้นตอนนี้เกิดบนชิปแสดงผลเฉพาะทาง ไม่กินทรัพยากร GPU ของ SoC
  \item \textbf{ขยายความละเอียด}: ยกภาพความละเอียดต่ำขึ้นเป็นความละเอียดสูงด้วยอัลกอริทึม AI ภาพคมชัดขึ้น
  \item \textbf{ลดการใช้พลังงาน}: เพราะการแทรกเฟรมและการขยายความละเอียดเป็นหน้าที่ของชิปแสดงผลเฉพาะทาง ภาระของ SoC จึงลดลง ความร้อนและการใช้ไฟก็ลดตาม เกมจึงเสถียรขึ้น
\end{itemize}

\booksubsection{การระบุตำแหน่ง GPS: ความถี่เดียวกับความถี่คู่}

โทรศัพท์มือถือระบุตำแหน่งหลักตามพึ่งพาGPSชั้นป้องกันดาวนำทางระบบ

\begin{itemize}
  \item \textbf{GPS ความถี่เดียว}: รับสัญญาณเฉพาะย่านความถี่ L1 ความแม่นยำราว 3-5 เมตร ในพื้นที่โล่งๆ ใช้ได้สบาย แต่ในเมืองที่ตึกสูงชิดกันเป็นแอ่งกระทะ การสะท้อนสัญญาณจะทำให้ตำแหน่งเบี่ยงเบนได้มาก
  \item \textbf{GPS ความถี่คู่}: รับสัญญาณทั้ง L1 และ L5 พร้อมกัน ย่าน L5 คลื่นยาวกว่า ทนการรบกวนได้ดีกว่า การคำนวณร่วมสองความถี่ช่วยตัดความคลาดเคลื่อนจากชั้นไอโอโนสเฟียร์ได้เกือบหมด ความแม่นยำถึงระดับ 1-3 เมตร ใช้นำทางหรือเรียกรถในเมือง จุดจะขยับน้อยลงและเสถียรขึ้นชัดเจน
  \item \textbf{รองรับหลายระบบดาวเทียม}: โทรศัพท์ยุคนี้ไม่ได้รองรับแค่ GPS ยังรองรับ BeiDou (จีน) GLONASS (รัสเซีย) Galileo (ยุโรป) และระบบดาวเทียมอีกหลายระบบ ยิ่งเชื่อมต่อดาวเทียมได้หลายดวงพร้อมกัน ยิ่งระบุตำแหน่งได้เร็วและแม่นยำ
\end{itemize}

\booksubsection{เซ็นเซอร์: เสาอากาศเล็กๆ ของโทรศัพท์}

ในโทรศัพท์มีเซ็นเซอร์มากกว่าสิบตัวทำงานเงียบๆ:

\begin{itemize}
  \item \textbf{เซ็นเซอร์วัดความเร่ง + ไจโรสโคป}: รับรู้ท่าทางของโทรศัพท์ เอาไว้หมุนมุมมองเวลาเล่นเกม และหมุนหน้าจออัตโนมัติ
  \item \textbf{เซ็นเซอร์แสง}: ปรับความสว่างหน้าจออัตโนมัติ
  \item \textbf{เซ็นเซอร์ระยะ}: ตอนโทรศัพท์แล้วแนบใกล้หู จอจะดับเองอัตโนมัติ
  \item \textbf{เซ็นเซอร์ลายนิ้วมือ}: ตอนนี้กระแสหลักคือลายนิ้วมือใต้หน้าจอแบบออปติคอลหรืออัลตราโซนิก
  \item \textbf{การรู้จำใบหน้า}: Face ID ของแอปเปิลใช้โครงสร้างแสง 3D ปลอดภัยระดับสูง ส่วนแอนดรอยด์จำนวนมากใช้การจับภาพ 2D ความปลอดภัยระดับทั่วไป
  \item \textbf{NFC}: การสื่อสารระยะใกล้ แตะเพื่อใช้บัตรรถเมล์ บัตรเข้าออกอาคาร บัตรธนาคาร
  \item \textbf{บารอมิเตอร์}: วัดความสูงจากระดับน้ำทะเล แม้ขึ้นบันไดก็ตรวจจับได้
  \item \textbf{เซ็นเซอร์ฮอลล์}: ตรวจจับการเปิดปิดของเคสแม่เหล็ก
\end{itemize}

\booksection{ระบบปฏิบัติการ: Android VS iOS VS HarmonyOS}

ฮาร์ดแวร์คือร่างกาย ระบบปฏิบัติการคือวิญญาณ ทั้งสามระบบมีเอกลักษณ์เฉพาะตัว:

\begin{itemize}
  \item \textbf{แอนดรอยด์}: เปรียบเหมือนชุมชนโอเพนซอร์ส ผู้ผลิตโทรศัพท์ปรับแต่งได้ลึก ทางเลือกเยอะ คุณเปลี่ยนธีมได้ ติดตั้งสโตร์แอปภายนอกได้ ลองเล่นได้หลากหลาย แต่ก็เพราะเปิดกว้าง ปัญหาความแตกฉานจึงรุนแรง ประสบการณ์ต่างแบรนด์ต่างกันมาก
  \item \textbf{iOS}: เปรียบเหมือนสวนที่ถูกจัดการอย่างประณีต แอปเปิลควบคุมทั้งซอฟต์แวร์และฮาร์ดแวร์ด้วยมือตัวเอง ประสบการณ์จึงเป็นเอกภาพและลื่นไหล คุณภาพแอปสูง ความเป็นส่วนตัวได้รับการคุ้มครองดี แต่ความอิสระต่ำ หลายเรื่องต้องทำตามกติกาของแอปเปิล
  \item \textbf{ฮาร์โมนีโอเอส}: สิ่งมีชีวิตชนิดใหม่ของหัวเว่ย ขายจุดเด่นที่ความสามารถแบบกระจายศูนย์ ทำให้อุปกรณ์ต่างชนิดทำงานร่วมกันได้ไร้รอยต่อ เดี๋ยวเราเจาะลึกกันภายหลัง
\end{itemize}

ระบบนิเวศของทั้งสามระบบล้วนสมบูรณ์ร่ำรวยมาก เลือกตามไลฟ์สไตล์ที่คุณชอบก็ดีแล้ว ไม่มีดีหรือเลวแบบสัมบูรณ์ มีแต่เหมาะหรือไม่เหมาะเท่านั้น

\booksection{วิเคราะห์พอร์ตเชื่อมต่อโทรศัพท์อย่างครบถ้วน}

พอร์ตโทรศัพท์เป็นสะพานเชื่อมต่อกับโลกภายนอก พอร์ตต่างชนิดกันมีความแตกต่างมาก:

\booksubsection{Micro USB}

\begin{itemize}
  \item \textbf{ลักษณะทางกายภาพ}: หัวเสียบรูปสี่เหลี่ยมคางหมูแบน แยกด้านหน้าหลัง (เสียบผิดข้างก็เข้าไม่เข้า) ตอนนี้ในโทรศัพท์เลิกใช้ไปเกือบหมดแล้ว แต่อุปกรณ์ระดับล่างบางตัวและของเก่ายังใช้อยู่
  \item \textbf{ข้อจำกัดการถ่ายโอน}: รองรับเฉพาะมาตรฐาน USB 2.0 อัตราถ่ายโอนสูงสุด 480Mbps กำลังชาร์จมักไม่เกิน 10W
  \item \textbf{ข้อเสีย}: หน้าหลังใช้ร่วมกันไม่ได้ อายุการเสียบถอดสั้น ถ่ายโอนช้า กำลังชาร์จต่ำ
\end{itemize}

\booksubsection{Lightning (หัวเสียบของแอปเปิล ตอนนี้เลิกใช้แล้ว)}

\begin{itemize}
  \item \textbf{ลักษณะทางกายภาพ}: หัวเสียบ 8 พินที่แอปเปิลพัฒนาเอง เสียบได้ทั้งสองด้าน ใช้กับ iPhone 14 และรุ่นก่อนหน้า ตอนนี้เลิกใช้ไปหมดแล้ว
  \item \textbf{ข้อจำกัดการถ่ายโอน}: หัวเสียบ Lightning ทางกายภาพรองรับเฉพาะมาตรฐาน USB 2.0 อัตราถ่ายโอนสูงสุด 480Mbps แม้ใช้สาย USB-C ต่อ Lightning ความเร็วในการถ่ายโอนก็ยังถูกจำกัดด้วยตัวหัวเสียบ Lightning เอง
  \item \textbf{สาเหตุที่เลิกใช้}: สหภาพยุโรปผ่านกฎหมายบังคับให้ตั้งแต่ปี 2024 โทรศัพท์ทุกเครื่องต้องใช้หัวเสียบ USB-C ร่วมกัน แอปเปิลจึงเปลี่ยนมาใช้ Type-C ทั้งหมดตั้งแต่ซีรีส์ iPhone 15 หัวเสียบ Lightning ก็อำลาเวทีประวัติศาสตร์ไปเป็นที่เรียบร้อย
  \item \textbf{สังเกต}: สายที่ได้การรับรอง MFi ของแอปเปิลจะมีชิปรับรองในตัว สายที่ไม่ผ่านการรับรองอาจชาร์จหรือถ่ายโอนข้อมูลได้ไม่ปกติ ถ้าคุณยังใช้ iPhone รุ่นเก่าอยู่ แนะนำให้เตรียมสาย Lightning เผื่อไว้
\end{itemize}

\booksubsection{USB Type-C}

\begin{itemize}
  \item \textbf{ลักษณะทางกายภาพ}: หัวเสียบรูปวงรี 24 พิน เสียบได้ทั้งสองด้าน ตอนนี้กลายเป็นหัวเสียบมาตรฐานของโทรศัพท์และอุปกรณ์อิเล็กทรอนิกส์ส่วนใหญ่แล้ว สหภาพยุโรปออกกฎบังคับตั้งแต่ปี 2024 ให้โทรศัพท์ แท็บเล็ต หูฟัง และอุปกรณ์พกพาใช้หัวเสียบ USB-C ร่วมกัน ภูมิภาคอื่นทั่วโลกก็กำลังตามรอย
  \item \textbf{โปรโตคอลการถ่ายโอน}: Type-C เป็นเพียงรูปแบบหัวเสียบ ความสามารถในการถ่ายโอนจริงขึ้นอยู่กับโปรโตคอลภายใน รองรับ USB 2.0 (480Mbps), USB 3.0/3.1/3.2 (5-20Gbps) หรือแม้แต่ USB4 (40Gbps)
  \item \textbf{การแยกแยะที่สำคัญ}:ไม่ใช่พอร์ต Type-C ทุกพอร์ตจะรองรับ USB 3.0! แม้โทรศัพท์หลายรุ่นจะใช้พอร์ต Type-C แต่ภายในเชื่อมต่อเพียงพินของ USB 2.0 ทำให้ความเร็วในการถ่ายโอนช้าเท่ากับพอร์ตแบบเก่า iPhone 16/17 รุ่นมาตรฐานรองรับเฉพาะความเร็ว USB 2.0 เท่านั้น ส่วนรุ่น Pro ถึงจะรองรับ USB 3.0 (ต้องใช้สาย USB-C to USB-C)
  \item \textbf{ความสามารถในการชาร์จ}: รองรับการชาร์จเร็วผ่านโปรโตคอลอย่าง USB PD สูงสุดได้ถึง 240W
  \item \textbf{เอาต์พุตวิดีโอ}: รองรับ DisplayPort Alt Mode ส่งภาพออกไปยังจอภาพผ่านสายต่อได้
\end{itemize}

\booksubsection{หัวเสียบเสียง 3.5mm}

\begin{itemize}
  \item \textbf{ลักษณะทางกายภาพ}: หัวเสียบกลมแบบ 3 ขาหรือ 4 ขา แบบ 3 ขาส่งเสียงซ้ายขวา แบบ 4 ขาเพิ่มช่องไมโครโฟนอีกหนึ่งทาง
  \item \textbf{คุณสมบัติ}: ส่งสัญญาณแบบอนาล็อก เสียบปุ๊บใช้ปั๊บ ไม่ต้องมีไดรเวอร์หรือตัวแปลง คุณภาพเสียงนิ่งคงที่ ดีเลย์เป็นศูนย์
  \item \textbf{สถานะปัจจุบัน}: โทรศัพท์ฟลากชิปส่วนใหญ่ตัดหัวเสียบ 3.5mm ออกไปแล้ว ต้องใช้สายต่อ Type-C หรือหูฟังบลูทูธ บางรุ่นระดับกลางและโทรศัพท์เกมมิ่งยังเก็บไว้
\end{itemize}

\booksection{ระบบผลิตภัณฑ์โทรศัพท์และกฎการตั้งชื่อ}

เผชิญกับสายผลิตภัณฑ์มากมายของแต่ละแบรนด์ การเข้าใจกฎการตั้งชื่อจะช่วยให้คุณหาโทรศัพท์ที่เหมาะกับตัวเองได้รวดเร็ว

\booksubsection{ซีรีส์ผลิตภัณฑ์ของแบรนด์หลัก:}

\begin{itemize}
  \item \textbf{Apple}:iPhone (รุ่นมาตรฐาน) / iPhone Plus (รุ่นหน้าจอใหญ่) / iPhone Pro (รุ่นสำหรับมืออาชีพ) / iPhone Pro Max (รุ่นหน้าจอใหญ่สำหรับมืออาชีพ) / iPhone SE (รุ่นราคาประหยัด) / iPhone 16e (รุ่นราคาประหยัดเน้น AI มาพร้อมชิป A18 และโมเด็ม C1 ที่ Apple พัฒนาเอง) ในปี 2026 ซีรีส์ iPhone 17 ทุกรุ่นจะใช้ชิป A19 Pro เต็มตัว ซึ่งจะยกระดับความสามารถด้าน AI ให้ดียิ่งขึ้นไปอีก
  \item \textbf{Huawei}: ซีรีส์ Mate (ฟลากชิปธุรกิจ) / ซีรีส์ P/Pura (ฟลากชิปกล้อง) / ซีรีส์ Mate X (ฟลากชิปพับ) / ซีรีส์ nova (ระดับกลางสายหนุ่มสาว) / ซีรีส์ Enjoy (ระดับเริ่มต้น)
  \item \textbf{Xiaomi}: Xiaomi รุ่นตัวเลข (ฟลากชิป) / Xiaomi MIX (ฟลากชิปสายสำรวจ/พับได้) / Redmi ซีรีส์ K (ฟลากชิปสายเพอร์ฟอร์แมนซ์) / Redmi ซีรีส์ Note (ระดับกลาง) / Redmi รุ่นตัวเลข (ระดับเริ่มต้น)
  \item \textbf{OPPO}: ซีรีส์ Find X (ฟลากชิป) / ซีรีส์ Reno (ระดับกลาง) / ซีรีส์ A (ระดับเริ่มต้น) / OnePlus รุ่นตัวเลข (ฟลากชิปสายเพอร์ฟอร์แมนซ์)
  \item \textbf{vivo}: ซีรีส์ X (ฟลากชิปกล้อง) / iQOO รุ่นตัวเลข (ฟลากชิปสายเพอร์ฟอร์แมนซ์) / ซีรีส์ S (ระดับกลางเน้นเซลฟี่) / ซีรีส์ Y (ระดับเริ่มต้น)
  \item \textbf{Samsung}: Galaxy S ซีรีส์ (แฟลกชิป) / Galaxy Z Fold/Flip (แฟลกชิปจอพับ) / Galaxy A ซีรีส์ (ระดับกลาง)
\end{itemize}

\booksubsection{ความหมายของคำต่อท้ายที่พบบ่อย:}

\begin{itemize}
  \item \textbf{Pro}: ย่อมาจาก Professional แปลว่ารุ่นมืออาชีพ มักอัปเกรดกล้อง หน้าจอ ชิป และสเปกหลักจากรุ่นมาตรฐาน
  \item \textbf{Max}: รุ่นใหญ่สุด หน้าจอใหญ่สุด แบตใหญ่สุด แต่สเปกไม่จำเป็นต้องสูงกว่า Pro
  \item \textbf{Ultra}: รุ่นเหนือระดับ มักเป็นสเปกสูงสุดของซีรีส์นั้น ใส่ของเต็มระบบ ราคาก็แพงสุด เช่น Samsung S24 Ultra, Xiaomi 14 Ultra
  \item \textbf{SE}: ย่อมาจาก Special Edition รุ่นพิเศษ ในทางปฏิบัติส่วนใหญ่คือรุ่นราคาเข้าถึงง่าย สเปกถูกลดลง ราคาถูกลง เช่น iPhone SE
  \item \textbf{Plus}: รุ่นขยายใหญ่ หน้าจอใหญ่กว่ารุ่นมาตรฐาน แบตก็ใหญ่ขึ้น แต่สเปกโดยรวมเหมือนเดิม
  \item \textbf{Lite/รุ่นเยาวชน}: รุ่นเบา สเปกถูกตัดทอน ราคาถูกลง
  \item \textbf{Note}: เดิมหมายถึงโทรศัพท์จอใหญ่ ตอนนี้แต่ละแบรนด์ใช้ไม่เหมือนกัน Redmi Note เป็นซีรีส์ระดับกลาง
  \item \textbf{Neo}: รุ่นสายใหม่คมกริบ มักหมายถึงรุ่นระดับกลางที่เน้นเพอร์ฟอร์แมนซ์
\end{itemize}

\booksection{เบสแบนด์และเสาอากาศ: ฮีโร่หลังฉากของการสื่อสารโทรศัพท์}

โทรศัพท์โทรและเข้าอินเทอร์เน็ตได้ด้วยคู่เบสแบนด์ชิปและเสาอากาศ

\booksubsection{ชิปเบสแบนด์:}

เบสแบนด์คือตัวประมวลผลที่รับผิดชอบการสื่อสารไร้สายในโทรศัพท์ มันมอดูเลตสัญญาณดิจิทัลเป็นคลื่นวิทยุส่งออกไป และดีมอดูเลตคลื่นวิทยุที่รับมากลับเป็นสัญญาณดิจิทัล ประสิทธิภาพของเบสแบนด์ส่งผลโดยตรงต่อคุณภาพการโทร ความเร็วอินเทอร์เน็ต และความนิ่งของสัญญาณ Qualcomm เป็นผู้จำหน่ายเบสแบนด์โทรศัพท์รายใหญ่ที่สุด MediaTek และซัมซุงก็มีเบสแบนด์พัฒนาเอง แอปเปิลเริ่มติดตั้งชิปเบสแบนด์ C1 ที่พัฒนาเองตั้งแต่ iPhone 16e เป็นต้นมา ถือเป็นจุดที่แอปเปิลเข้าสู่ยุคเบสแบนด์พัฒนาเองอย่างเป็นทางการ C1 ใช้กระบวนการผลิต 4nm ประหยัดไฟเป็นเลิศ ส่วน iPhone 16/16 Plus/16 Pro/16 Pro Max ยังใช้เบสแบนด์ Qualcomm (Snapdragon X75 5G) คาดว่ารุ่นต่อๆ ไปจะค่อยๆ ย้ายมาใช้เบสแบนด์ของตัวเอง ย้อนไปก่อนหน้านั้น iPhone ซีรีส์ XS/11 ที่ใช้เบสแบนด์ Intel สัญญาณไม่ค่อยดี เป็นจุดปวดหัวของผู้ใช้ไม่น้อย

\booksubsection{การออกแบบเสาอากาศ:}

โทรศัพท์มือถือเสาอากาศซ่อนที่ตัวเครื่องในเครื่อง, โดยทั่วไปล้อมพันขอบผ้าตั้ง เสาอากาศออกแบบตรงเชื่อมต่อมีผลสัญญาณรับประสิทธิภาพ:

\begin{itemize}
  \item \textbf{เทคโนโลยี MIMO}: รับหลายทางส่งหลายทาง ใช้เสาอากาศหลายตัวรับส่งข้อมูลพร้อมกัน เพิ่มความเร็วและความนิ่ง 4×4 MIMO เร็วกว่า 2×2 MIMO
  \item \textbf{จำนวนเสาอากาศ}: โทรศัพท์ฟลากชิปมักมีเสาอากาศสิบกว่าตัว (เซลลูลาร์ Wi-Fi บลูทูธ NFC GPS ต้องมีเสาอากาศแยกของใครของมัน)
  \item \textbf{ผลกระทบจากตัวเครื่องโลหะ}: โลหะจะบังสัญญาณ โทรศัพท์ตัวเครื่องโลหะจึงต้องออกแบบแถบเสาอากาศ (แถบพลาสติกบนขอบ) ส่วนตัวเครื่องแก้วและพลาสติกกระทบสัญญาณน้อยกว่า
\end{itemize}

\booksubsection{วิวัฒนาการเทคโนโลยีเครือข่าย:}

\begin{itemize}
  \item \textbf{4G LTE}: อัตราความเร็วสูงสุดทางทฤษฎี ดาวน์โหลด 1Gbps / อัปโหลด 500Mbps (Cat.18) ประสบการณ์จริงมัก 50-150Mbps ครอบคลุมกว้างที่สุด ใช้งานได้สบาย
  \item \textbf{5G NR}: อัตราความเร็วสูงสุดทางทฤษฎี ดาวน์โหลด 10Gbps+ ย่าน Sub-6GHz ครอบคลุมได้ดี mmWave มิลลิเมตรเวฟเร็วสุดขั้วแต่ครอบคลุมน้อยมาก
  \item \textbf{5G-A(5G Advanced)}:เป็นเวอร์ชันเสริมของ 5G หรืออาจเรียกว่า 5.5G มีอัตราความเร็วในการดาวน์โหลดสูงสุดถึง 10Gbps และอัปโหลดที่ 1Gbps พร้อมทั้งเพิ่มความหนาแน่นของการเชื่อมต่อและความแม่นยำในการระบุตำแหน่งอย่างมาก ตั้งแต่ปี 2024 เป็นต้นมา China Mobile, China Telecom และ China Unicom เริ่มปรับใช้เครือข่าย 5G-A ในเมืองสำคัญ โดยโทรศัพท์เรือธงอย่าง Huawei Mate 70/80 ซีรีส์, Xiaomi 15/16 ซีรีส์ ฯลฯ ได้รองรับ 5G-A แล้ว 5G-A ยังนำมาซึ่งความสามารถ "การรวมการสื่อสารและการรับรู้" ที่สถานีฐานสามารถสื่อสารและรับรู้สภาพแวดล้อมไปพร้อมกันได้ ซึ่งถือเป็นโครงสร้างพื้นฐานสำหรับเศรษฐกิจพื้นที่ต่ำ (การจัดส่งด้วยโดรน) และระบบเครือข่ายยานยนต์
\end{itemize}

\booksubsection{การแบ่งประเภทย่านความถี่เครือข่าย:}

ย่านความถี่เครือข่ายที่โทรศัพท์รองรับเป็นตัวกำหนดว่ามันใช้ได้กับประเทศไหน เครือข่ายของผู้ให้บริการไหน:

\begin{itemize}
  \item \textbf{FDD-LTE Band(Bซีรีส์)}: เช่น B1, B3, B5 เป็นย่านความถี่ 4G แต่ละประเทศใช้ย่านต่างกัน จีนใช้หลักๆ B1/B3/B5/B8/B34/B38/B39/B40/B41
  \item \textbf{TDD-LTE Band}: อีกหนึ่งมาตรฐานของ 4G จีนใช้ B38/B39/B40/B41 เป็นจำนวนมาก
  \item \textbf{5G NR Band(Nซีรีส์)}: เช่น n1, n3, n28, n41, n78, n79 n78 (3.5GHz) เป็นย่านหลักของ 5G จีน n41 (2.5GHz) ใช้ที่ China Mobile เยอะ
  \item \textbf{คำแนะนำก่อนซื้อ}: ถ้าต้องเอาไปใช้ต่างประเทศ ตรวจสอบก่อนซื้อว่าโทรศัพท์รองรับย่านความถี่ของประเทศปลายทางหรือไม่ โทรศัพท์ที่วางขายในจีนมักรองรับย่านของจีนครบที่สุด แต่อาจขาดย่านต่างประเทศบางส่วน
\end{itemize}

\booksection{โทรศัพท์ของคุณเล่นบทบาทกี่อย่าง?}

มันไม่ใช่แค่โทรศัพท์อีกต่อไป:

\begin{enumerate}
  \item \textbf{เครื่องมือสื่อสาร}: โทรศัพท์ ข้อความ SMS WeChat อีเมล……
  \item \textbf{กล้องอัจฉริยะ}: ถ่ายรูป อัดวิดีโอ สแกนเอกสาร……
  \item \textbf{กระเป๋าเงิน}: จ่ายเงิน บัตรรถเมล์ บัตรธนาคาร บัตรสมาชิก……
  \item \textbf{ประตูสู่ข้อมูล}: ข่าว สภาพอากาศ หุ้น ความรู้……
  \item \textbf{เครื่องเล่นเกม}: ตั้งแต่เกมจับคู่จนถึงเกม 3A ระดับทุ่มเท……
  \item \textbf{สำนักงานเคลื่อนที่}: เอกสาร ตาราง ประชุม ทำงานร่วมกัน……
  \item \textbf{ผู้ช่วยด้านสุขภาพ}: นับก้าว วัดชีพจร ติดตามการนอน ออกกำลังกาย……
  \item \textbf{รีโมทสมาร์ทโฮม}: ควบคุมไฟ แอร์ ทีวี หุ่นยนต์ดูดฝุ่น……
  \item \textbf{บัตรประจำตัวดิจิทัล}: บัตรประชาชนอิเล็กทรอนิกส์ ใบขับขี่อิเล็กทรอนิกส์ รหัสสุขภาพ……
  \item \textbf{ศูนย์ความบันเทิง}: เพลง วิดีโอ หนังสืออิเล็กทรอนิกส์ พอดแคสต์……
\end{enumerate}

ลองคิดดู อุปกรณ์เล็กๆ เครื่องเดียว รองรับชีวิตได้มากแค่ไหน! ดังนั้น การเรียนรู้ที่จะใช้มันให้ดี คือการเรียนรู้ที่จะใช้ชีวิตให้ดีขึ้น

\begin{tcbpractice}{ลงมือปฏิบัติ: ทำความเข้าใจการกำหนดค่าฮาร์ดแวร์}
ใน [ตั้งค่า] ดูว่าหาโมเดลโปรเซสเซอร์ RAM และความจุ ROM ได้ไหม ทำความเข้าใจความสามารถของโทรศัพท์คุณ ผู้ใช้แอนดรอยด์สามารถดาวน์โหลด “AnTuTu” หรือ “DevCheck” เพื่อดูข้อมูลฮาร์ดแวร์โดยละเอียด

\end{tcbpractice}
